\documentclass[12pt,a4paper,oneside]{article}
\usepackage[brazil]{babel}
\usepackage[utf8]{inputenc}
\usepackage{olymp}

\setcounter{problem}{3}

\begin{document}

\begin{problem}{Palíndromo}{palin}{2}
 
Palíndromo é uma palavra ou frase que permanece a mesma quando lida de trás pra frente. Ou seja, é a mesma quando lida tanto da esquerda para direita quanto da direita para esquerda.

Exemplos: ARARA, OSSO, ACAIACA, RADAR.

Existem frases palíndromas, como ``Roma me tem amor'' e a famosa ``Socorram-me, subi no ônibus em Marrocos!''. Neste exercício você não deve se preocupar com frases, apenas com palavras.


%\Notes
%
%Observações, se tiver.

\InputFile

A entrada contém apenas uma linha com a palavra a ser testada. Restrições: a palavra contém no máximo 30 letras, todas elas maiúsculas.

\OutputFile

Seu programa deve gerar uma linha de saída para cada caso de teste, contendo a palavra ``\texttt{PALINDROMO}'' ou ``\texttt{NAO}'', em maiúsculas e sem acentos, indicando se a palavra da entrada é palíndroma ou não.

% \Notes
% Preste atenção aos tipos das variáveis, pois $F(90) \approx 2^{61}$.

\Example

\begin{example}
\exmp{
ARARA
}{
PALINDROMO
}%
\end{example}

\begin{example}
\exmp{
PALINDROMO
}{
NAO
}%
\end{example}

\begin{example}
\exmp{
BANANA
}{
NAO
}%
\end{example}

\begin{example}
\exmp{
ABCDEFFEDCBA
}{
PALINDROMO
}%
\end{example}

%Comentários sobre o exemplo, se necessário.

\end{problem}

\end{document}
